% chapters/02_apriori_representations.tex
%
% Part I, Chapter 2.  A-priori representations: Fourier, Wavelet,
% Scattering.  Adapted from the geometric tour's Chapter 2 with the
% Scattering Transform significantly expanded (it's the bridge to
% Chapter 3) and the connection to convolution made more explicit.

\chapter{A-priori representations}
\labch{apriori}

\epigraph{Music is the pleasure the human mind experiences from counting
without being aware that it is counting.}{Gottfried Wilhelm Leibniz}

%------------------------------------------------------------------------
\section{The question and three answers}
%------------------------------------------------------------------------

\noindent A sustained violin note arriving at our ears is a pressure
wave: a long sequence of small fluctuations around atmospheric pressure,
sampled by the eardrum at something like 40\,000 times per second.  As a
mathematical object it is a vector in $\R^N$ with $N$ very large.  As a
listening experience, it is a single sustained tone with a particular
timbre, attack, and decay.

The gap between those two descriptions is what a \emph{representation} is
supposed to close.  We want a transformation $\Phi$ such that $\Phi(\xb)$
exposes the structure we care about --- pitch, brightness, attack, the
texture of the bowing --- and ideally hides what we don't.  The previous
chapter argued that any reasonable $\Phi$ looks like an inner product
with a chosen set of reference vectors.  This chapter looks at three
specific choices, each historically and practically important.  All
three are \emph{a-priori} in the sense that the reference vectors are
designed mathematically, not learned from data.  The next chapter will
turn to representations whose reference vectors come from data.

The three a-priori representations we consider are:

\begin{itemize}
    \item The \emph{Discrete Fourier Transform} (Section~\ref{sec:dft}),
    whose reference vectors are pure sinusoids.  Fourier exposes
    frequency content with maximum sharpness and is the workhorse of
    classical signal processing.

    \item The \emph{Discrete Wavelet Transform} (Section~\ref{sec:dwt}),
    whose reference vectors are localized oscillations --- a single
    \emph{mother wavelet}, scaled and translated to cover the
    time-frequency plane.  Wavelets give us frequency information with
    time localisation, at the cost of some spectral precision.

    \item The \emph{Scattering Transform} (Section~\ref{sec:scattering}),
    a multi-layer cascade of wavelet convolutions and absolute-value
    nonlinearities.  It is the bridge between classical signal processing
    and modern convolutional neural networks: structurally identical to a
    CNN, but with all the filters fixed by mathematics rather than learned
    from data.
\end{itemize}

\noindent The progression Fourier $\to$ Wavelet $\to$ Scattering will
turn out to add, one step at a time, three things: \emph{locality}
(wavelets are localised, sinusoids are not), \emph{depth} (Scattering
stacks two or three wavelet transforms), and \emph{nonlinearity}
(Scattering inserts an absolute value between the layers).  Once we have
these three, a fourth step --- replacing the predefined filters with
\emph{learned} ones --- gets us a convolutional network.  That fourth
step is the subject of the next chapter.

%------------------------------------------------------------------------
\section{The Discrete Fourier Transform}
\label{sec:dft}
%------------------------------------------------------------------------

The Discrete Fourier Transform represents a length-$N$ signal as a sum
of $N$ complex sinusoids.  For each frequency index $k \in
\{0, 1, \ldots, N-1\}$,
\begin{equation}
    \hat{\xb}[k]
    \;=\;
    \sum_{n=0}^{N-1} \xb[n] \, e^{-i 2\pi k n / N}.
    \label{eq:dft}
\end{equation}
The inverse transform reconstructs $\xb$ from $\hat\xb$ via the same
formula with the opposite sign in the exponent and a $1/N$ normalisation:
\begin{equation}
    \xb[n]
    \;=\;
    \frac{1}{N} \sum_{k=0}^{N-1} \hat{\xb}[k] \, e^{i 2 \pi k n / N}.
    \label{eq:idft}
\end{equation}

\marginnote{%
The DFT in \eqref{eq:dft} would naively cost $O(N^2)$ operations to
evaluate.  The \emph{Fast Fourier Transform} reduces this to
$O(N \log N)$ by exploiting the recursive structure of the sum.  The FFT
is what makes Fourier analysis practical for audio; we use it
extensively in Part~II.}

\paragraph{The DFT as projection.}  Compare \eqref{eq:dft} with the
inner-product formula \eqref{eq:inner} from the previous chapter.  The
DFT is just the inner product of $\xb$ with the $k$-th complex sinusoid:
\begin{equation}
    \hat{\xb}[k] \;=\; \inner{\xb}{\mathbf{f}_k},
    \qquad
    \mathbf{f}_k[n] \;=\; e^{i 2 \pi k n / N}.
    \label{eq:dft_as_inner}
\end{equation}
Each $\mathbf{f}_k$ is one of our reference vectors: a discrete sinusoid
that completes exactly $k$ full cycles over the length of the signal.
The set $\{\mathbf{f}_k\}_{k=0}^{N-1}$ is an \emph{orthogonal basis} for
$\C^N$, and in fact an orthonormal one after the $1/N$ scaling in
\eqref{eq:idft}.  This is why the inverse transform exists in closed
form: an orthogonal basis is its own decoder.

In matrix form, letting $\mathbf{F}$ be the $N \times N$ matrix with
$\mathbf{F}_{kn} = e^{-i 2 \pi k n / N}$,
\begin{equation}
    \hat\xb \;=\; \mathbf{F}\, \xb,
    \qquad
    \xb \;=\; \frac{1}{N} \mathbf{F}^H\, \hat\xb.
    \label{eq:dft_matrix}
\end{equation}
The Fourier transform is, exactly, a particular choice of
$\mathbf{B}$ in the framework of the previous chapter.

\paragraph{What the DFT exposes.}  Because each basis vector
$\mathbf{f}_k$ is a pure tone of frequency $k$, the coefficient
$\hat\xb[k]$ is large precisely when $\xb$ contains energy at that
frequency.  The magnitude $\abs{\hat\xb[k]}$ is the amount of frequency
$k$ in the signal; the argument $\arg \hat\xb[k]$ is its phase.  For
musical signals --- a sustained pitched tone --- the DFT magnitude
shows a peak at the fundamental and at each harmonic, telling us
immediately how loud each harmonic is.  This is the basis of pitch
detection, spectral envelopes, and timbre analysis.

\paragraph{What the DFT hides.}  Two properties of \eqref{eq:dft_as_inner}
matter.  First, the sum runs over the \emph{entire} signal: every sample
contributes to every coefficient.  A short click at $n = 100$
contributes to $\hat\xb[k]$ for every $k$.  The DFT has no notion of
``when'' --- it gives a global summary of frequency content, averaged
over all of time.  Second, the sinusoids $\mathbf{f}_k$ extend
unattenuated across the whole signal length; they are global functions
masquerading as basis vectors.

This is fine for stationary signals (an organ note holding constant for
seconds at a time) but disastrous for transient signals (a snare hit,
a consonant in speech, the attack of a piano note).  For such signals
we want \emph{locality}: a representation that tells us not just what
frequencies are present but \emph{when}.

\paragraph{Translation invariance and frequency sensitivity.}
The DFT has one beautiful invariance and one painful sensitivity.

If we shift $\xb$ in time by $\Delta$, the DFT magnitudes are
unchanged --- only the phases change:
\begin{equation}
    \xb[n - \Delta] \;\;\longleftrightarrow\;\;
    \hat\xb[k] \, e^{-i 2 \pi k \Delta / N}.
\end{equation}
This is what we want: a violin note arriving slightly later sounds the
same.  The magnitude spectrum is \emph{translation-invariant}.

But if we shift $\xb$ in \emph{frequency} by even a small amount --- say,
the violin is slightly mistuned, or our sample rate drifts ---
the DFT coefficients can move around dramatically.  A pure tone at
frequency $k + 0.5$ smears across many bins because it doesn't match
any single basis vector.  This is the well-known \emph{leakage}
phenomenon, and it's the fundamental cost of using strictly periodic
basis vectors on signals that are not strictly periodic.

\marginnote{%
A continuous version of the leakage problem is the
\emph{Heisenberg uncertainty principle for signals}: a representation
cannot have arbitrarily sharp resolution in both time and frequency.
The DFT chooses maximum frequency resolution (sharp basis vectors in
frequency) and accepts zero time resolution (global basis vectors in
time).  Wavelets choose a different trade-off; the scattering transform
yet another.}

\paragraph{Convolution becomes multiplication.}  The Fourier transform's
most useful algebraic property, far beyond its role as a representation,
is that convolution in the time domain corresponds to pointwise
multiplication in the frequency domain:
\begin{equation}
    \yb \;=\; \xb \conv \mathbf{h}
    \qquad \Longleftrightarrow \qquad
    \hat\yb[k] \;=\; \hat\xb[k] \cdot \hat{\mathbf{h}}[k].
    \label{eq:conv_theorem}
\end{equation}
This is the \emph{convolution theorem}; we sketch a proof in
Appendix~\ref{ch:convolution_theorem}.  The practical consequence is
that a convolution that would naively cost $O(NM)$ to evaluate directly
costs $O(N \log N)$ via three FFTs (forward, multiply, inverse).  For
large kernels, the FFT-based convolution is by far the faster method.

In the other direction, multiplication of two signals corresponds to
\emph{convolution} of their spectra.  This is how the Short-Time
Fourier Transform works (next), why amplitude modulation in audio
creates sidebands, and why the spectral analysis of a windowed signal
shows the window's spectrum convolved with the signal's spectrum.

%------------------------------------------------------------------------
\section{The Discrete Wavelet Transform}
\label{sec:dwt}
%------------------------------------------------------------------------

The DFT gives us frequency at the cost of time.  The Wavelet Transform
restores time by giving up some frequency.

\paragraph{Mother wavelets and their family.}  Pick a \emph{mother
wavelet} $\psi$: a short, oscillatory, zero-mean function of finite
energy.  Generate a family of basis vectors by scaling and translating
it:
\begin{equation}
    \psi_{a, b}[n]
    \;=\;
    \frac{1}{\sqrt{|a|}} \, \psi\!\left( \frac{n - b}{a} \right).
    \label{eq:wavelet_family}
\end{equation}
The scale parameter $a$ stretches or compresses the wavelet: a large
$a$ produces a wide, low-frequency wavelet; a small $a$ produces a
narrow, high-frequency one.  The translation parameter $b$ shifts the
wavelet along the time axis.  The normalisation $1/\sqrt{|a|}$ keeps the
energy constant across scales.

The Discrete Wavelet Transform of a signal $\xb$ is the array of inner
products of $\xb$ with this family:
\begin{equation}
    \mathbf{w}[a, b] \;=\; \inner{\xb}{\psi_{a, b}}.
\end{equation}
The coefficient $\mathbf{w}[a, b]$ tells us how much of the signal looks
like a wavelet of scale $a$ centred at time $b$.  Large coefficients
mark places in time-frequency where the signal has structure at that
particular scale.

\marginnote{%
Common mother wavelets include the \emph{Haar} (a square wave: simplest,
discontinuous), \emph{Daubechies} families (compactly supported,
orthogonal, smoother), the \emph{Mexican hat} (a Gaussian's second
derivative), and the \emph{Morlet} (a complex sinusoid times a Gaussian,
much used in time-frequency audio analysis).  The choice of mother
wavelet matters in practice; for this book the Morlet is the most
relevant.}

\paragraph{The time-frequency plane.}  Where the DFT gives one
coefficient per frequency, the DWT gives one coefficient per
\emph{location in time-frequency}.  Plot a DWT as a heatmap with time
on the horizontal axis and scale on the vertical axis, and you see a
picture: which frequencies are active when.  This is exactly the
\emph{spectrogram-like} visualisation that audio engineers use to look
at sounds.  In Part~II we build a related construction --- the Short-Time
Fourier Transform --- which gives a similar picture using DFTs of
short windowed segments.  Wavelets are the more principled relative,
with proper mathematical foundations rather than ad-hoc windowing.

\paragraph{What the wavelet representation costs.}  We gain time
localisation, but we pay for it.  Two specific costs.

First, the representation is no longer translation-invariant in
magnitude.  Shifting $\xb$ in time \emph{does} change which wavelet
coefficients are large, because the coefficients are indexed by both
scale and position.  A wavelet representation is translation-\emph{equivariant}:
the pattern of coefficients shifts in lockstep with the signal, but the
shifted pattern is not the same pattern.

Second, the frequency resolution at any single scale is lower than the
DFT's.  Each wavelet covers a range of frequencies rather than a single
one; the trade is wider time coverage in exchange for narrower
frequency coverage.  This is unavoidable: tight in time forces broad in
frequency, and vice versa.

\paragraph{Stability to frequency drift.}  In exchange, the DWT is much
more stable than the DFT to small drifts in frequency.  A pure tone at
frequency $f_0 + \delta$ produces wavelet coefficients close to those
of a pure tone at $f_0$, because each wavelet's frequency response is
already broad.  This makes wavelets the better tool for analysing
real-world signals whose frequencies wobble or evolve --- vibrato,
glissando, the slowly drifting pitch of a struck bell.

\paragraph{Matrix form.}  Just as with the DFT, the DWT has a matrix
representation.  Let $\boldsymbol\Psi$ be the matrix whose rows are the
wavelets $\psi_{a, b}$ at every $(a, b)$ pair in our chosen
discretisation.  Then
\begin{equation}
    \mathbf{w} \;=\; \boldsymbol\Psi\, \xb,
    \qquad
    \xb \;=\; \boldsymbol\Psi^H\, \mathbf{w}
    \quad\text{(if the wavelets form an orthonormal basis).}
    \label{eq:dwt_matrix}
\end{equation}
We have the same analysis-synthesis structure as the DFT, with a
different choice of basis.  The framework is invariant; the basis is
the design choice.

%------------------------------------------------------------------------
\section{Between transforms: the time-frequency trade}
%------------------------------------------------------------------------

Before we get to the Scattering Transform, it is worth pausing on the
relationship between the DFT and the DWT, because it makes the design
space visible.

The DFT picks one extreme of the time-frequency trade: maximally sharp
in frequency, maximally blurry in time.  Every basis vector spans the
whole signal, so we know exactly which frequencies are present but
nothing about when.

The Short-Time Fourier Transform (STFT), which we build in Part~II,
picks a middle point: we cut the signal into short overlapping
windows, take the DFT of each, and stack the results.  Each window has
fixed time width (call it $W$), so every frequency bin has the same
time resolution $W$ and the same frequency resolution $\sim\!1/W$.
Audio engineers reach for the STFT constantly; it is the basis of the
spectrogram.

The DWT picks a different middle point, one where the time-frequency
resolution varies with scale.  High frequencies (small $a$) get fine
time resolution and coarse frequency resolution --- a clear ``the
transient happened here'' answer.  Low frequencies (large $a$) get
coarse time resolution and fine frequency resolution --- a clear
``the pitch is $f_0$'' answer over a longer window.  This is exactly
the right behaviour for natural signals: short events tend to be
broadband, long events tend to be narrowband.

\marginnote{%
The DWT's variable resolution is sometimes called \emph{constant-Q
analysis}: each band has a fixed ratio of frequency width to centre
frequency.  This matches the way the human auditory system processes
sound, where critical bands also have approximately constant $Q$.  It is
not a coincidence that wavelet-like decompositions appear in the
cochlea.}

The progression is from rigid (DFT) to adaptive (DWT) to multi-layered
(Scattering, next).  Each step gives up some elegance --- the DFT has a
closed-form inverse; the DWT requires careful choice of mother wavelet;
the Scattering Transform has multiple stages and a fixed but non-trivial
schedule of operations --- in exchange for better behaviour on real
signals.

%------------------------------------------------------------------------
\section{The Scattering Transform}
\label{sec:scattering}
%------------------------------------------------------------------------

The Scattering Transform, introduced by Mallat in the early 2010s
(\emph{Group invariant scattering}, Communications on Pure and Applied
Mathematics, 2012), builds on the DWT in two crucial ways: it adds a
nonlinearity (the absolute value) between successive wavelet
decompositions, and it stacks several of them.  The resulting structure
is, mathematically and conceptually, a convolutional neural network with
all the filters fixed by hand.

\paragraph{One layer.}  We start with a signal $\xb$ and a family of
wavelets $\{\psi_{\lambda_1}\}$ indexed by scale (we collapse the
notation $\psi_{a, b}$ into $\psi_\lambda$ where $\lambda$ ranges over
both scale and shift).  The first-layer scattering coefficients are
\begin{equation}
    S_1 \xb \;=\; \abs{\xb \conv \psi_{\lambda_1}}.
    \label{eq:scattering1}
\end{equation}
Compute the wavelet transform of $\xb$, take the absolute value of each
coefficient.  This is the only operation; nothing learned, nothing
adjustable.

\paragraph{Two layers.}  The output $S_1 \xb$ is itself a signal --- in
fact, a collection of signals, one per scale.  We can apply the same
operation to each of them, using a second family of wavelets
$\{\psi_{\lambda_2}\}$:
\begin{equation}
    S_2 \xb \;=\; \abs{\, \abs{\xb \conv \psi_{\lambda_1}} \conv \psi_{\lambda_2} \,}.
    \label{eq:scattering2}
\end{equation}
This is a wavelet transform applied to the magnitude of a wavelet
transform.  The result captures \emph{interactions} between frequency
bands: how does activity at scale $\lambda_1$ vary over time?  How does
that variation itself look at scale $\lambda_2$?

The transform can be extended to three or more layers by repeating the
recipe.  In practice, two or three layers suffice; later layers tend
to capture diminishing amounts of energy.

\paragraph{Why the absolute value.}  Without the nonlinearity between
the two wavelet transforms, equation \eqref{eq:scattering2} would
collapse into a single wavelet transform (because the composition of
two linear operators is a linear operator).  The absolute value
prevents this collapse: it is the \emph{nonlinearity} that makes a deep
representation more expressive than a shallow one.

The choice of absolute value, specifically, is not arbitrary.  It is
the simplest nonlinearity that is \emph{stable} (Lipschitz: small
changes in input produce only small changes in output) and that
removes the phase information (which is highly oscillatory in time and
would otherwise destabilise the second layer).  Other nonlinearities
work too, but the absolute value has the cleanest theoretical
properties.

\paragraph{Stability and invariance.}  The scattering transform inherits
two key properties from its construction:

\begin{itemize}
    \item It is \emph{stable to deformations}: if $\xb$ is replaced by
    a slightly time-warped version $\xb \circ \tau$ where $\tau$ is a
    smooth deformation, the scattering coefficients $S\xb$ change by an
    amount controlled by the smoothness of $\tau$, not by the size of
    the deformation itself.  This is the right notion of robustness for
    natural signals, which are constantly subject to small distortions
    (recording-room acoustics, microphone characteristics, breath
    fluctuations).

    \item It is \emph{translation-invariant at large scale}: when the
    final layer is averaged over a window much wider than the signal's
    features, the result depends on $\xb$'s content but not on where in
    time that content occurs.  This is what we want for classification
    (a violin is a violin whether it arrives at sample $1000$ or
    $2000$).
\end{itemize}

\marginnote{%
\textbf{Why this matters.}  Bruna and Mallat showed in 2013 (\emph{Invariant
scattering convolution networks}, IEEE TPAMI) that a two-layer scattering
transform achieves near-state-of-the-art accuracy on some image and audio
classification tasks, \emph{without any learning}.  The features it
produces are good enough that a simple linear classifier on top reaches
accuracies that a CNN of the same depth would achieve.  This is a
remarkable result: it suggests that much of what early CNN layers do can
be replicated by fixed mathematical operations.  Of course, deeper
learned networks eventually outperform scattering, but the comparison
frames the question: what is the depth of learning truly buying us?}

\paragraph{The architecture.}  Look at the structure of the scattering
transform from a distance.  Each layer convolves the input with a bank
of filters, applies a nonlinearity, and passes the result to the next
layer.  Strip the words ``wavelet'' and ``absolute value'' and replace
them with ``learned filter'' and ``ReLU'' and you have, exactly, a
convolutional neural network.  Wavelet convolution becomes
$\Wb \conv \xb$; the absolute value becomes ReLU; the cascade of
layers is the same.

This is the bridge from classical signal processing to modern deep
learning.  The scattering transform is the \emph{minimum} deep network:
filters chosen for theoretical optimality (translation invariance,
locality, stability), nonlinearity chosen for stability.  A CNN
\emph{replaces} the principled-but-rigid wavelets with filters that the
data adjusts, trading mathematical guarantees for flexibility.  Some
problems do not benefit from that trade --- and for those, scattering is
competitive.  Most problems do --- and for those, we need learning.

\paragraph{Hierarchical features.}  The scattering transform also makes
explicit a structural point about \emph{depth}.  The first-layer
coefficients $\abs{\xb \conv \psi_{\lambda_1}}$ are very similar to a
spectrogram: they tell us about frequency content over time.  The
second-layer coefficients capture how that frequency content is itself
modulated.  Energy that fluctuates rapidly at the first layer becomes
visible as a clear peak at a particular $\lambda_2$ in the second layer.
This is the modulation spectrum, central in audio analysis.

We will see the same hierarchical structure in CNNs: first layer
detects edges; second layer detects combinations of edges; third layer
detects parts; eventually the network detects objects.  The CNN's
hierarchy is learned; the scattering transform's is mathematical.  They
are, structurally, the same hierarchy.

%------------------------------------------------------------------------
\section{Closing thoughts}
%------------------------------------------------------------------------

We have walked through three a-priori representations: the DFT
(maximum frequency resolution, no time resolution), the DWT (a
flexible time-frequency trade-off), and the Scattering Transform
(multi-layer, locally invariant, deformation-stable).  Each is a
specific choice of $\Phi(\xb) = f(\mathbf{B}\xb)$ in the framework of
the previous chapter, with $\mathbf{B}$ a fixed matrix designed by
mathematicians and $f$ either the identity, the magnitude
$\abs{\cdot}$, or a composition of those.

The three transforms answer different questions about a signal.  The
DFT answers ``what frequencies are present?''  The DWT answers ``what
frequencies are present, and when?''  The Scattering Transform
answers ``what frequencies are present, when, and how do they
interact?''  Each successive question is more demanding; each requires
more structure in the representation.  The progression suggests a
natural fourth step: ``what features are present, and we don't even
know what features we mean?''  When we don't know what kinds of
patterns to look for, we let the data tell us.  That is the subject of
the next chapter --- representations whose reference vectors are
\emph{learned}.

The progression also previews the architectural pattern we will
re-encounter many times: \emph{projection}, \emph{nonlinearity},
\emph{stacking}.  These three together make depth meaningful.  Without
projection there is nothing to compose; without nonlinearity composition
collapses; without stacking there is no hierarchy.  The Scattering
Transform demonstrates that all three can be in place without any
learning at all.  The point of learning, as we will see, is to let the
projections themselves adapt to whatever signals we throw at the
network --- to choose $\mathbf{B}$ from data rather than from a
mathematician's foresight.
